\section{Conclusión}

La contribución principal de este artículo es metodológica: establecemos un benchmark unificado, reproducible y justo para métodos de interpolación EEG. Al someter todos los algoritmos---incluyendo líneas base clásicas---a optimización rigurosa de hiperparámetros mediante Optuna, demostramos que la brecha de rendimiento entre métodos tradicionales y modernos es sustancialmente menor de lo típicamente reportado en la literatura. Esto no es una debilidad sino una fortaleza: aclara que el rendimiento realista requiere ajuste cuidadoso, no solo novedad algorítmica.

Dentro de este marco de evaluación justo, surge un hallazgo empírico claro: los métodos de regularización temporal, particularmente TRSS, logran calidad de reconstrucción competitiva con MAE promedio de \ResTVMAE\ y SNR de \ResTVSNR~dB, superando líneas base geométricas exhaustivamente optimizadas como Spherical Spline (MAE $\approx$ \ResBaselineMAE, SNR $\approx$ \ResBaselineSNR~dB).

Más allá de métricas de error estándar, la regularización TRSS y Tikhonov demuestran superioridad robusta en fidelidad fisiológica. TRSS preserva dinámicas de señal (DTW $\approx$ \ResTVDTW), minimiza distorsión de fase y recupera con precisión Potenciales Relacionados con Eventos (ERP) y Densidad Espectral de Potencia (PSD) incluso bajo pérdida severa de canales contiguos (p. ej., 40\% de canales faltantes cercanos). Esto respalda la conclusión de que incorporar hipótesis de suavidad temporal en el framework del Laplaciano de grafo espacial proporciona una ventaja sobre modelos puramente espaciales.

\subsection*{Validación Estadística a Nivel de Prueba}
Un análisis pairwise exhaustivo a nivel de prueba en 736 iteraciones y 4,336 muestras agregadas confirma la robustez estadística de estas conclusiones. Utilizando la prueba Mann-Whitney U con tamaño de efecto delta de Cliff, realizamos 55 comparaciones de métodos pareados después de excluir el valor atípico basado en visibilidad del ranking. Los resultados demuestran significancia estadística fuerte: 39 comparaciones alcanzaron $p < 0.05$, y 38 llegaron a $p < 0.001$. Los métodos espacial-temporal y temporal (temporal\_laplacian, TRSS, IDW, linear) se agruparon estrechamente en el nivel superior (MAE mediano: 0.075--0.081) sin separación estadísticamente significativa. Esta confirmación a nivel de prueba valida los hallazgos a nivel de iteración y establece una jerarquía de rendimiento clara y robusta en contextos experimentales diversos.

La validación extensiva en conjuntos de datos de PhysioNet, BCI Competition IV 2a y MNE Sample confirma la generalización de estos hallazgos. Este trabajo proporciona un punto de referencia para investigación futura en interpolación EEG, enfatizando la importancia de evaluación justa y optimización rigurosa de hiperparámetros. El trabajo futuro debe explorar si el rendimiento en tiempo real puede ser alcanzado mediante aceleración GPU o estrategias de aproximación, y si la sintonización específica de tarea o sujeto produce ganancias de rendimiento adicionales.
